% =====================================================================
\section{Analiza problemu}
\label{sec:analiza}

\subsection{Pomysł na aplikację}

\textbf{PlayLink} to lekka platforma webowa typu \emph{LFG} (\textit{Looking-For-Group})
służąca do natychmiastowego znajdowania współgraczy do gier niszowych, modów oraz
tytułów retro. Projekt odpowiada na trzy konkretne problemy:

\begin{itemize}
  \item \textbf{Martwe społeczności} — niszowe i stare gry mają zbyt mało aktywnych
        graczy, by samodzielnie utrzymać żywe lobby.
  \item \textbf{Brak dedykowanych narzędzi} — nie istnieje proste narzędzie do szybkiego
        zestawiania partnerów do takich gier.
  \item \textbf{Wysokie bariery wejścia} — istniejące platformy wymagają rejestracji,
        są skomplikowane lub nieaktywne.
\end{itemize}

Odpowiedzią PlayLink jest natychmiastowy dostęp \emph{bez klasycznej rejestracji}
(tożsamość BIP39), przeglądarka aktywnych sesji w czasie rzeczywistym oraz prosty
mechanizm udostępniania pokoju linkiem.

\subsection{Analiza istniejących rozwiązań}
\label{sec:konkurencja}

Przeanalizowano trzy najpopularniejsze alternatywy. Dla każdej wskazano co najmniej
dwie zalety i dwie wady wraz z uzasadnieniem, dlaczego dana cecha jest istotna w
kontekście gier niszowych. Zrzuty ekranu omawianych aplikacji znajdują się w
prezentacji projektowej (załącznik).

\paragraph{Discord — kanały LFG na serwerach.}
Tekstowe kanały „looking-for-group” na serwerach społeczności konkretnych gier, z
komunikacją głosową i tekstową w jednym miejscu.

\begin{itemize}
  \item \textbf{Zalety:}
  \begin{itemize}
    \item Bardzo szybka i bezpośrednia komunikacja (głos + tekst) — kluczowa przy
          umawianiu rozgrywki na żywo.
    \item Ogromna, aktywna baza użytkowników — łatwo znaleźć ludzi do gier popularnych.
  \end{itemize}
  \item \textbf{Wady:}
  \begin{itemize}
    \item Brak dedykowanego systemu sesji — ogłoszenia toną w strumieniu wiadomości i
          szybko znikają, co uniemożliwia ich przeglądanie.
    \item Trzeba znać i dołączyć do konkretnego serwera; brak wyszukiwania ofert
          \emph{między} serwerami — przy grach niszowych właściwe kanały są martwe.
  \end{itemize}
\end{itemize}

\paragraph{Reddit — r/LFG i subreddity gier.}
Tablice ogłoszeń „szukam grupy” w formacie postów na tematycznych subredditach.

\begin{itemize}
  \item \textbf{Zalety:}
  \begin{itemize}
    \item Duża baza użytkowników i prosty, znany format ogłoszeń.
    \item Asynchroniczność — ogłoszenie można zostawić i wrócić do odpowiedzi później.
  \end{itemize}
  \item \textbf{Wady:}
  \begin{itemize}
    \item Całkowity brak komunikacji w czasie rzeczywistym — czeka się na komentarz,
          a posty starzeją się w godzinach, nie w minutach.
    \item Brak filtrowania po grze, regionie czy poziomie oraz brak struktury sesji;
          subreddity gier niszowych są często martwe.
  \end{itemize}
\end{itemize}

\paragraph{FACEIT — matchmaking e-sportowy.}
Profesjonalna platforma e-sportowa z matchmakingiem, ligami i anti-cheatem dla gier
kompetytywnych (CS2, Dota~2, League of Legends).

\begin{itemize}
  \item \textbf{Zalety:}
  \begin{itemize}
    \item Zaawansowany dobór graczy na podstawie umiejętności (skill-based matchmaking).
    \item Duża baza graczy skupionych wokół rywalizacji oraz dedykowane serwery.
  \end{itemize}
  \item \textbf{Wady:}
  \begin{itemize}
    \item Obsługuje wyłącznie popularne tytuły e-sportowe — całkowite pominięcie gier
          niszowych i retro.
    \item Wymaga instalacji klienta i anti-cheata, a część funkcji jest płatna —
          zbyt wysoka bariera dla casualowego gracza.
  \end{itemize}
\end{itemize}

\paragraph{Podsumowanie.}
Żadne z istniejących rozwiązań nie łączy jednocześnie: (1) braku barier wejścia,
(2) struktury sesji z aktualizacją w czasie rzeczywistym oraz (3) wsparcia dla gier
niszowych. PlayLink przejmuje najlepsze elementy konkurencji (czat sesyjny jak na
Discordzie, otwartą listę ogłoszeń jak na Reddicie), eliminuje ich wady (trwała,
przeszukiwalna lista pokojów z filtrowaniem; aktualizacja przez WebSocket) i dodaje
nowe cechy: logowanie bezhasłowe BIP39 oraz kryptograficznie podpisywane wiadomości.

\subsection{Wymagania funkcjonalne}
\label{sec:wf}

Zidentyfikowano 15 wymagań funkcjonalnych; wszystkie zostały zrealizowane w prototypie.

\renewcommand{\arraystretch}{1.25}
\begin{longtable}{@{}clc@{}}
\toprule
\textbf{ID} & \textbf{Wymaganie} & \textbf{Status} \\
\midrule
\endhead
WF-01 & Rejestracja i logowanie bezhasłowe (fraza BIP39 + podpis ECDSA) & \done \\
WF-02 & Odzyskiwanie konta przez import frazy mnemonicznej & \done \\
WF-03 & Edycja nazwy użytkownika (z filtrem wulgaryzmów) & \done \\
WF-04 & Tworzenie i konfiguracja pokoju (gra, limit graczy, lokalizacja) & \done \\
WF-05 & Metadane pokoju (opis, link do komunikatora, wymagania) & \done \\
WF-06 & Przeglądanie listy aktywnych pokojów na żywo & \done \\
WF-07 & Dołączanie do pokoju i jego opuszczanie & \done \\
WF-08 & Czat tekstowy w obrębie pokoju (w czasie rzeczywistym) & \done \\
WF-09 & Kryptograficzne podpisywanie wiadomości czatu (EIP-191) & \done \\
WF-10 & Wyszukiwanie pokojów i użytkowników & \done \\
WF-11 & Filtrowanie pokojów po grze oraz sortowanie listy & \done \\
WF-12 & Kalendarz pokoju i RSVP dla zaplanowanej sesji & \done \\
WF-13 & Własne gry (custom games) i ping zależny od lokalizacji lobby & \done \\
WF-14 & Panel administratora: zarządzanie pokojami i kategoriami gier & \done \\
WF-15 & Moderacja: wyrzucenie uczestnika z pokoju (kick) przez admina & \done \\
\bottomrule
\end{longtable}

\subsection{Wymagania niefunkcjonalne i ich miary}
\label{sec:wnf}

\renewcommand{\arraystretch}{1.3}
\begin{longtable}{@{}cp{0.30\textwidth}p{0.46\textwidth}c@{}}
\toprule
\textbf{ID} & \textbf{Wymaganie} & \textbf{Miara / kryterium akceptacji} & \textbf{St.} \\
\midrule
\endhead
WN-01 & Wydajność & Czas odpowiedzi API $<$ 200~ms dla 95\% żądań (p95) & \done \\
WN-02 & Responsywność realtime & Aktualizacja UI po zdarzeniu $<$ 1000~ms & \done \\
WN-03 & Dostępność & Uptime $\geq$ 99\% w skali miesiąca & \done \\
WN-04 & Skalowalność & Do 128 użytkowników jednocześnie w pokoju & \done \\
WN-05 & Niezawodność & Docker healthcheck + automatyczny restart kontenera & \done \\
WN-06 & Realtime / odporność & WebSocket z heartbeatem 30~s i auto-reconnectem & \done \\
WN-07 & Bezpieczeństwo transportu & TLS~1.3 dla HTTPS i WSS (Let's Encrypt) & \done \\
WN-08 & Ochrona przed nadużyciami & Rate limiting: 10/min (REST), 10~wiad./10~s (czat WS) & \done \\
WN-09 & Użyteczność & Utworzenie pokoju w $\leq$ 30~s bez instrukcji & \done \\
WN-10 & Modułowość & Frontend i backend wdrażane niezależnie & \done \\
WN-11 & Jakość kodu & Pokrycie testami backendu $\geq$ 80\% (bramka CI) & \done \\
WN-12 & Wdrażalność & Czas od \code{push} do wersji live $\approx$ 3~min & \done \\
\bottomrule
\end{longtable}

\subsection{Przypadki użycia i scenariusze}
\label{sec:uc}

W systemie występują czterej aktorzy: \textbf{Gość} (niezalogowany), \textbf{Użytkownik}
(zalogowany, dziedziczy uprawnienia Gościa), \textbf{Administrator} (dziedziczy
uprawnienia Użytkownika) oraz \textbf{System} (zadania automatyczne — weryfikacja
challenge-response i czyszczenie wygasłych pokojów).

\paragraph{Lista przypadków użycia.}
\begin{itemize}
  \item \textbf{Gość:} przeglądanie listy pokojów; wyszukiwanie i filtrowanie;
        rejestracja / logowanie.
  \item \textbf{Użytkownik:} tworzenie i konfiguracja pokoju; dołączanie/opuszczanie;
        czat (podpisany); planowanie wydarzenia i RSVP; edycja nazwy; odzyskiwanie konta.
  \item \textbf{Administrator:} wyrzucenie uczestnika (kick); zarządzanie kategoriami gier.
  \item \textbf{System:} weryfikacja podpisu i wydanie JWT; czyszczenie wygasłych
        pokojów i wiadomości.
\end{itemize}

Diagram przypadków użycia przedstawiono na rys.~\ref{fig:uc}. Poniżej opisano
scenariusze pięciu kluczowych przypadków użycia — po jednym opracowanym przez
każdego członka zespołu — w ujęciu: przepływ główny oraz przepływ alternatywny.

\paragraph{UC-1: Rejestracja i pierwsze logowanie (BIP39). \normalfont\textit{(autor: Nikodem Nowak)}}
\begin{itemize}
  \item \textit{Aktor:} Gość. \quad \textit{Warunek wstępny:} brak.
  \item \textit{Przepływ główny:}
  \begin{enumerate}
    \item Gość wybiera „Enter realm” i generuje nową 12-słowną frazę BIP39
          (klucz prywatny derywowany lokalnie, nigdy nie opuszcza przeglądarki).
    \item Klient pobiera z backendu jednorazowy \emph{nonce} (\code{POST /auth/request-nonce}).
    \item Klient podpisuje komunikat zawierający nonce (\code{personal_sign}, EIP-191).
    \item Backend odtwarza adres z podpisu, weryfikuje nonce i wydaje token JWT
          (\code{POST /auth/verify}).
    \item BFF zapisuje JWT w ciasteczku \code{httpOnly}; użytkownik jest zalogowany.
  \end{enumerate}
  \item \textit{Przepływ alternatywny:} nonce wygasł lub został już użyty —
        backend odrzuca weryfikację (HTTP 400/401), klient ponawia od kroku~2.
\end{itemize}

\paragraph{UC-2: Utworzenie pokoju i rozpoczęcie czatu. \normalfont\textit{(autor: Bartosz Kołaciński)}}
\begin{itemize}
  \item \textit{Aktor:} Użytkownik. \quad \textit{Warunek wstępny:} zalogowany; mniej niż
        3 aktywne pokoje.
  \item \textit{Przepływ główny:}
  \begin{enumerate}
    \item Użytkownik wypełnia formularz (nazwa, gra, limit graczy, lokalizacja,
          opcjonalnie opis/link/wymagania) i zatwierdza.
    \item Backend tworzy pokój, dołącza twórcę jako członka i rozgłasza nową listę
          pokojów przez kanał \code{/ws/rooms}.
    \item Użytkownik wchodzi do pokoju; otwiera się kanał czatu
          \code{/ws/rooms/<name>/chat}.
  \end{enumerate}
  \item \textit{Przepływ alternatywny:} przekroczony limit 3 pokojów lub zajęta nazwa —
        backend zwraca błąd walidacji, formularz wyświetla komunikat.
\end{itemize}

\paragraph{UC-3: Dołączenie do pokoju i deklaracja RSVP. \normalfont\textit{(autor: Jakub Rosiak)}}
\begin{itemize}
  \item \textit{Aktor:} Użytkownik. \quad \textit{Warunek wstępny:} pokój istnieje i nie
        jest pełny.
  \item \textit{Przepływ główny:}
  \begin{enumerate}
    \item Użytkownik wybiera pokój z listy i dołącza; roster jest rozgłaszany do
          pozostałych członków.
    \item Jeśli pokój ma zaplanowane wydarzenie, użytkownik deklaruje obecność
          (\emph{present} / \emph{maybe} / \emph{absent}); zdarzenie RSVP jest rozgłaszane.
    \item Użytkownik bierze udział w czacie pokoju.
  \end{enumerate}
  \item \textit{Przepływ alternatywny:} pokój pełny lub wygasł — dołączenie zostaje
        odrzucone, użytkownik wraca do listy.
\end{itemize}

\paragraph{UC-4: Moderacja — wyrzucenie uczestnika. \normalfont\textit{(autor: Jakub Rusek)}}
\begin{itemize}
  \item \textit{Aktor:} Administrator. \quad \textit{Warunek wstępny:} adres administratora
        znajduje się w \code{ADMIN_ADDRESSES}.
  \item \textit{Przepływ główny:}
  \begin{enumerate}
    \item Administrator wybiera uczestnika i akcję „kick”.
    \item Backend weryfikuje uprawnienia, usuwa członkostwo i rozgłasza ramkę
          \emph{kick} oraz zaktualizowany roster; wyrzucony klient zostaje rozłączony.
  \end{enumerate}
  \item \textit{Przepływ alternatywny:} żądanie pochodzi od konta bez uprawnień
        administratora — backend odrzuca operację (HTTP 403).
\end{itemize}

\paragraph{UC-5: Wyszukiwanie i filtrowanie pokojów. \normalfont\textit{(autor: Mateusz Kosowski)}}
\begin{itemize}
  \item \textit{Aktor:} Gość lub Użytkownik. \quad \textit{Warunek wstępny:} brak.
  \item \textit{Przepływ główny:}
  \begin{enumerate}
    \item Aktor otwiera listę pokojów, która jest na bieżąco aktualizowana przez
          kanał \code{/ws/rooms}.
    \item Aktor wpisuje frazę wyszukiwania i/lub wybiera filtr gry oraz porządek
          sortowania.
    \item Lista zawęża się do pokojów spełniających kryteria; wybór pokoju prowadzi
          do jego szczegółów (a po zalogowaniu — do dołączenia, UC-2/UC-3).
  \end{enumerate}
  \item \textit{Przepływ alternatywny:} brak pokojów spełniających kryteria —
        wyświetlany jest stan pusty z podpowiedzią rozszerzenia filtrów lub utworzenia
        własnego pokoju.
\end{itemize}

\subsection{Diagram przypadków użycia}
\label{sec:uc-diagram}

\begin{figure}[H]
  \centering
  \includegraphics[width=0.95\textwidth]{uc.png}
  \caption{Diagram przypadków użycia PlayLink (aktorzy: Gość, Użytkownik,
  Administrator, System).}
  \label{fig:uc}
\end{figure}

\subsection{Wstępna architektura rozwiązania}
\label{sec:architektura}

PlayLink ma architekturę trójwarstwową z wzorcem \textbf{BFF}
(\textit{Backend-for-Frontend}). Przeglądarka komunikuje się \emph{bezpośrednio} z
API (REST oraz dwa kanały WebSocket) dla większości odczytów i całego ruchu realtime,
a z serwerem SvelteKit (BFF) tylko dla operacji wymagających dostępu do ciasteczka
sesyjnego \code{httpOnly} (logowanie/wylogowanie oraz uwierzytelnione akcje formularzy),
które BFF proxy-uje do backendu z dołączonym tokenem JWT.

\begin{itemize}
  \item \textbf{Wzorce projektowe:} BFF (izolacja JWT od JavaScriptu klienta),
        wstrzykiwanie zależności (FastAPI \code{Depends}), wzorzec repozytorium przez
        SQLModel/SQLAlchemy, publish–subscribe dla rozgłaszania zdarzeń WebSocket.
  \item \textbf{Główne moduły:} klient Svelte (podpisywanie BIP39 przez ethers/viem),
        BFF SvelteKit (adapter-node), API FastAPI (REST + WebSocket + zadanie w tle),
        baza PostgreSQL.
  \item \textbf{Stos technologiczny} zestawiono w tab.~\ref{tab:stack}.
\end{itemize}

\renewcommand{\arraystretch}{1.2}
\begin{table}[H]
\centering\small
\begin{tabularx}{\textwidth}{@{}lY@{}}
\toprule
\textbf{Warstwa} & \textbf{Technologia} \\
\midrule
Frontend & SvelteKit~2 / Svelte~5 (runes), TypeScript, Tailwind CSS (build: Bun) \\
Podpisywanie & ethers + viem (BIP39 + EIP-191) \\
Backend & FastAPI (Python 3.14+), serwer ASGI (build: uv) \\
ORM / migracje & SQLModel (SQLAlchemy + Pydantic), Alembic \\
Baza danych & PostgreSQL~18 (prod) / SQLite in-memory (testy) \\
Realtime & FastAPI WebSockets (kanały: rooms, chat) \\
Uwierzytelnianie & tożsamość BIP39 + ECDSA (\code{eth_account}), JWT HS256 \\
Rate limiting & slowapi (REST) + przesuwne okno per-adres (czat) \\
Orkiestracja & Docker Compose (db, backend, frontend), Nginx, Tailscale \\
CI/CD & GitHub Actions $\rightarrow$ deploy SSH przez Tailscale \\
\bottomrule
\end{tabularx}
\caption{Stos technologiczny PlayLink.}
\label{tab:stack}
\end{table}

Przepływ uwierzytelniania challenge-response przedstawia rys.~\ref{fig:auth}.

\begin{figure}[H]
  \centering
  \includegraphics[width=0.92\textwidth]{sekwencja-auth.png}
  \caption{Diagram sekwencji: bezhasłowe logowanie (challenge-response BIP39/ECDSA).
  Klucz prywatny nigdy nie opuszcza przeglądarki; JWT trafia do ciasteczka
  \texttt{httpOnly} ustawianego przez BFF.}
  \label{fig:auth}
\end{figure}

\subsection{Diagram komponentów}
\label{sec:komponenty}

\begin{figure}[H]
  \centering
  \includegraphics[width=0.85\textwidth]{komponenty.png}
  \caption{Diagram komponentów: przeglądarka, Nginx (TLS), BFF SvelteKit, API FastAPI
  z zadaniem czyszczącym oraz baza PostgreSQL; wdrożenie przez Tailscale.}
  \label{fig:komponenty}
\end{figure}

\subsection{Wstępny model rozwiązania: warstwa danych i warstwa logiki}
\label{sec:model}

\subsubsection*{Warstwa danych}

Trwałość zapewnia PostgreSQL z ośmioma encjami zarządzanymi przez SQLModel i migracjami
Alembic. Model danych przedstawia diagram ERD na rys.~\ref{fig:er}.

\begin{itemize}
  \item \code{user} — tożsamość: \code{identity_address} (adres EIP-55 odtworzony z ECDSA)
        oraz unikalna \code{username}; \emph{brak hasła}.
  \item \code{nonce} — jednorazowe, krótkożyciowe wyzwania; w bazie przechowywany jest
        wyłącznie skrót SHA-256 wartości nonce (\code{value}), z flagą \code{used}.
  \item \code{room} — sesja gry: nazwa (unikalna), gra, lokalizacja lobby, limit graczy,
        metadane oraz \code{expires_at} (TTL, domyślnie 60~min).
  \item \code{roommember} — tabela łącząca (M:N) użytkowników i pokoje (klucz złożony).
  \item \code{message} — wiadomości czatu; opcjonalny podpis EIP-191 (\code{signature})
        i podpisany znacznik czasu (\code{sent_at}) zapewniają weryfikowalne autorstwo.
  \item \code{roomevent} / \code{roomeventrsvp} — zaplanowane wydarzenie pokoju oraz
        deklaracje obecności (status \emph{present}/\emph{absent}/\emph{maybe},
        semantyka upsert).
  \item \code{game} — katalog kategorii gier (kolejność wg \code{sort_order}).
\end{itemize}

\begin{figure}[H]
  \centering
  \includegraphics[width=\textwidth]{er.png}
  \caption{Model danych (ERD) PlayLink — osiem encji wraz z kluczami i relacjami.}
  \label{fig:er}
\end{figure}

\subsubsection*{Warstwa logiki}

Logika backendu jest skupiona w usłudze FastAPI udostępniającej REST oraz dwa kanały
WebSocket. Najważniejsze elementy:

\begin{itemize}
  \item \textbf{REST} — endpointy uwierzytelniania (\code{/auth/request-nonce},
        \code{/auth/verify}), CRUD pokojów, profil użytkownika, administracja.
        Dokumentacja generowana automatycznie (Swagger/OpenAPI).
  \item \textbf{WebSocket} — kanał \code{/ws/rooms} (rozgłasza pełną listę pokojów po
        każdej mutacji) oraz \code{/ws/rooms/<name>/chat} (czat z autoryzacją JWT i
        kontrolą członkostwa; ramki: event, RSVP, roster, room-closed, kick).
  \item \textbf{Logika biznesowa} — m.in.: limit 3 aktywnych pokojów na użytkownika,
        weryfikacja podpisów EIP-191 z odrzuceniem nadmiernego dryfu czasu
        (\code{CHAT_SIGNATURE_SKEW_SECONDS}), upsert RSVP, dobór pingu na podstawie
        lokalizacji lobby (odległość Haversine).
  \item \textbf{Zadanie w tle} — cykliczne czyszczenie wygasłych pokojów i wiadomości
        (co \code{ROOM_CLEANUP_INTERVAL_SECONDS}).
\end{itemize}
